\section{Experiments}

\subsection{Full-grid coverage}

The completed Informer matrix contains \(972\) eligible seed-runs: 36 dataset-horizon cells, nine activation configurations, and three seeds per configuration. All \(324\) dataset-horizon-configuration cells have complete three-seed coverage under the matched batch-size-32, six-epoch protocol. \Cref{tab:phase4_coverage} summarizes the coverage split across ETTh, ETTm, and Solar.

\begin{table*}[t]
\centering
\caption{Phase-4 full-grid coverage. A main-table cell is eligible only when all three seeds are present under batch size 32 and six training epochs.}
\label{tab:phase4_coverage}
\resizebox{\textwidth}{!}{%
\begin{tabular}{lrrrr}
\toprule
Group & dataset-horizon cells & config cells & seed-runs & 3-seed config cells \\
\midrule
ETTh & 10 & 90 & 270/270 & 90/90 \\
ETTm & 10 & 90 & 270/270 & 90/90 \\
Solar & 16 & 144 & 432/432 & 144/144 \\
All & 36 & 324 & 972/972 & 324/324 \\
\bottomrule
\end{tabular}%
}
\end{table*}


The aggregation is strict. A row is reusable only when dataset, horizon, sequence length, label length, layer counts, ProbSparse factor, batch size, epochs, activation signature, and seed all match the canonical key. The Solar4 source file initially produced non-finite losses because six target values were missing; those values were repaired before relaunching the exact missing Solar4 seed-runs. The repaired source checksum is recorded in the artifact notes.

\subsection{Activation-only changes beat GELU in most cells}

\Cref{tab:phase4_best_activation} reports the best activation in each completed dataset-horizon cell. A non-GELU activation is best in \(34/36\) cells. The Wilson interval for this descriptive cell proportion is \(81.9\%\) to \(98.5\%\), but the more important point is the breadth of the pattern: all 20 ETT/ETTm cells prefer a non-GELU activation, and 14 of 16 Solar cells do as well. The only GELU-best cases are Solar3 at horizons 24 and 96.

\begin{table*}[t]
\centering
\caption{Best activation per dataset-horizon in the completed Phase-4 grid. Improvements are relative MSE reductions against the same-protocol GELU row with matched seeds; GELU-best cells have no relative-improvement entry.}
\label{tab:phase4_best_activation}
\resizebox{\textwidth}{!}{%
\begin{tabular}{llcccc}
\toprule
Dataset & H & best config & MSE & gain vs GELU & seed win rate \\
\midrule
ETTh1 & 24 & softsign & 0.532 $\pm$ 0.026 & +5.2\% & 66.7\% \\
ETTh1 & 48 & softsign & 0.646 $\pm$ 0.073 & +12.6\% & 100.0\% \\
ETTh1 & 168 & tanh & 0.912 $\pm$ 0.131 & +12.9\% & 66.7\% \\
ETTh1 & 336 & softsign & 1.219 $\pm$ 0.042 & +10.4\% & 100.0\% \\
ETTh1 & 720 & tanh+out & 1.328 $\pm$ 0.030 & +6.0\% & 100.0\% \\
ETTh2 & 24 & tanh\_sin & 0.516 $\pm$ 0.023 & +27.8\% & 100.0\% \\
ETTh2 & 48 & softsign & 1.606 $\pm$ 0.386 & +38.8\% & 100.0\% \\
ETTh2 & 168 & tanh\_sin & 2.963 $\pm$ 0.509 & +41.3\% & 100.0\% \\
ETTh2 & 336 & tanh+out & 2.054 $\pm$ 0.219 & +47.7\% & 100.0\% \\
ETTh2 & 720 & tanh+out & 3.298 $\pm$ 0.186 & +18.6\% & 100.0\% \\
ETTm1 & 24 & tanh & 0.314 $\pm$ 0.015 & +13.4\% & 100.0\% \\
ETTm1 & 48 & tanh & 0.453 $\pm$ 0.023 & +10.8\% & 100.0\% \\
ETTm1 & 96 & tanh\_sin & 0.559 $\pm$ 0.050 & +9.3\% & 100.0\% \\
ETTm1 & 288 & tanh & 0.802 $\pm$ 0.044 & +7.2\% & 66.7\% \\
ETTm1 & 672 & softsign & 0.922 $\pm$ 0.023 & +7.3\% & 66.7\% \\
ETTm2 & 24 & tanh & 0.685 $\pm$ 0.191 & +42.5\% & 100.0\% \\
ETTm2 & 48 & softsign & 0.309 $\pm$ 0.014 & +21.6\% & 100.0\% \\
ETTm2 & 96 & tanh\_sin & 1.335 $\pm$ 0.214 & +16.4\% & 100.0\% \\
ETTm2 & 288 & tanh+out & 2.355 $\pm$ 0.359 & +55.7\% & 100.0\% \\
ETTm2 & 672 & tanh+out & 2.851 $\pm$ 0.162 & +62.7\% & 100.0\% \\
Solar1 & 24 & Swish & 0.150 $\pm$ 0.000 & +1.7\% & 66.7\% \\
Solar1 & 96 & softsign & 0.224 $\pm$ 0.001 & +2.6\% & 100.0\% \\
Solar2 & 24 & tanh & 0.066 $\pm$ 0.003 & +6.4\% & 100.0\% \\
Solar2 & 96 & tanh & 0.108 $\pm$ 0.005 & +10.8\% & 100.0\% \\
Solar3 & 24 & GELU & 0.198 $\pm$ 0.009 & -- & -- \\
Solar3 & 96 & GELU & 0.349 $\pm$ 0.007 & -- & -- \\
Solar4 & 24 & tanh\_sin & 0.220 $\pm$ 0.003 & +1.9\% & 100.0\% \\
Solar4 & 96 & enc-tanh & 0.290 $\pm$ 0.005 & +1.9\% & 66.7\% \\
Solar5 & 24 & Swish & 0.109 $\pm$ 0.002 & +4.2\% & 100.0\% \\
Solar5 & 96 & enc-tanh & 0.180 $\pm$ 0.003 & +7.6\% & 100.0\% \\
Solar6 & 24 & tanh\_sin & 0.265 $\pm$ 0.002 & +1.1\% & 100.0\% \\
Solar6 & 96 & tanh\_sin & 0.358 $\pm$ 0.009 & +2.3\% & 100.0\% \\
Solar7 & 24 & ReLU & 3.281 $\pm$ 1.035 & +12.0\% & 66.7\% \\
Solar7 & 96 & enc-tanh & 4.970 $\pm$ 1.353 & +38.2\% & 100.0\% \\
Solar8 & 24 & tanh\_sin & 0.216 $\pm$ 0.003 & +0.8\% & 100.0\% \\
Solar8 & 96 & softsign & 0.306 $\pm$ 0.005 & +2.6\% & 100.0\% \\
\bottomrule
\end{tabular}%
}
\end{table*}


The largest best-cell gains occur in ETTh2 and ETTm2. ETTh2/336 and ETTm2/288 select hidden \texttt{tanh} with output \texttt{tanh}; ETTm2/672 selects the same output-bounded configuration and has the largest relative gain in the table. Solar gains are smaller and more site-specific. This is useful evidence against a trivial explanation that any non-GELU activation always helps: Solar3 keeps GELU, while Solar7 has high variance and a ReLU best cell at horizon 24.

\subsection{Bounded and tanh-family activations are strong static choices}

\Cref{tab:phase4_static_ranking} ranks configurations within each dataset-horizon cell and averages the ranks over all 36 cells. The strongest static choices are \texttt{tanh\_sin001\_all}, \texttt{tanh\_all}, and \texttt{softsign\_all}. They also have high top-3 coverage and positive-cell rates against GELU. This supports the main positive lesson: bounded sign-preserving nonlinearities are not merely one-off winners, but robust candidates under this Informer protocol.

\begin{table*}[t]
\centering
\caption{Static activation ranking across the 36 completed dataset-horizon cells. Rank is computed within each cell over the nine Phase-4 configurations.}
\label{tab:phase4_static_ranking}
\resizebox{\textwidth}{!}{%
\begin{tabular}{lrrrrr}
\toprule
Config & mean rank & best cells & top-3 cells & mean gain & positive cells \\
\midrule
tanh\_sin & 2.78 & 8 & 27 & +11.7\% & 83.3\% \\
tanh & 3.19 & 7 & 25 & +11.9\% & 80.6\% \\
softsign & 3.47 & 8 & 22 & +11.2\% & 80.6\% \\
enc-tanh & 4.83 & 3 & 11 & +5.1\% & 66.7\% \\
Swish & 4.89 & 2 & 5 & +3.6\% & 77.8\% \\
dec-tanh & 5.67 & 0 & 3 & +1.9\% & 66.7\% \\
GELU & 6.22 & 2 & 4 & +0.0\% & 0.0\% \\
ReLU & 6.83 & 1 & 4 & -2.3\% & 36.1\% \\
tanh+out & 7.11 & 5 & 7 & -36.3\% & 25.0\% \\
\bottomrule
\end{tabular}%
}
\end{table*}


\Cref{tab:phase4_static_uncertainty} adds the uncertainty view used to keep this claim scoped. The top three static configurations have positive mean gains with descriptive intervals that remain above zero when summarized over dataset-horizon cells. In contrast, decoder-only \texttt{tanh} is weak, ReLU is negative on average, and output \texttt{tanh} has a strongly negative mean despite winning five cells. Therefore the paper should not recommend output bounding as a default. It should be treated as a targeted placement candidate for selected long-horizon ETT/ETTm cells.

\begin{table}[t]
\centering
\caption{Cell-level uncertainty for static Informer activation choices. The mean gain is computed from paired-seed relative MSE improvements over same-protocol GELU, first averaged inside each dataset-horizon cell and then summarized over the 36 cells. Intervals are descriptive 95\% normal intervals over cells, not a claim of population-level optimality.}
\label{tab:phase4_static_uncertainty}
\resizebox{\textwidth}{!}{%
\begin{tabular}{lrrrr}
\toprule
Config & mean gain with 95\% interval & positive cells & best cells & top-3 cells \\
\midrule
\texttt{tanh\_all} & +11.9\% [ +6.6, +17.3 ] & 29/36 & 7 & 25 \\
\texttt{tanh\_sin001\_all} & +11.7\% [ +5.7, +17.7 ] & 30/36 & 8 & 27 \\
\texttt{softsign\_all} & +11.2\% [ +5.3, +17.0 ] & 29/36 & 8 & 22 \\
\texttt{enc\_tanh\_dec\_gelu} & +5.1\% [ +1.3, +9.0 ] & 24/36 & 3 & 11 \\
\texttt{swish\_all} & +3.6\% [ +0.5, +6.8 ] & 28/36 & 2 & 5 \\
\texttt{enc\_gelu\_dec\_tanh} & +1.9\% [ -0.7, +4.5 ] & 24/36 & 0 & 3 \\
\texttt{relu\_all} & -2.3\% [ -4.6, 0.0 ] & 13/36 & 1 & 4 \\
\texttt{tanh\_all\_outtanh} & -36.3\% [ -58.8, -13.9 ] & 9/36 & 5 & 7 \\
\bottomrule
\end{tabular}%
}
\end{table}


\subsection{Placement effects are real but conditional}

The best-cell counts are spread across configurations. \texttt{tanh\_sin001} and \texttt{softsign} each win eight cells, \texttt{tanh} wins seven, output \texttt{tanh} wins five, encoder-side \texttt{tanh} wins three, Swish wins two, GELU wins two, and ReLU wins one. Full-hidden bounded activations are the safest static candidates, but placement changes the result in important cells.

Output \texttt{tanh} is best in ETTh1/720, ETTh2/336, ETTh2/720, ETTm2/288, and ETTm2/672. These are high-gain cells, which explains why output bounding is tempting. However, the same placement has mean rank \(7.11\) and a negative average gain across the grid. Encoder-side \texttt{tanh} is best in Solar4/96, Solar5/96, and Solar7/96, while decoder-side \texttt{tanh} wins no cell. This pattern is consistent with placement being a conditional design variable rather than a universal trick.

\subsection{Regime factors are diagnostic, not a router}

Summary-level factor bins cover all 36 dataset-horizon cells and produce 783 factor-bin rows. They reveal substantial regime dependence: 196 factor-conditioned contrasts have absolute end-minus-start improvement at least \(0.02\). The strongest contrasts are often warnings rather than positive rules. For example, high volatility, jump intensity, and turbulence bins frequently show that output \texttt{tanh} becomes worse, not that a simple high-volatility rule should select it.

Sample-level leave-one-seed-out oracles are stricter. Among 84 leave-one-seed-out oracle rows, only seven are positive, all on Solar1/96. \Cref{tab:phase4_loso_oracle} lists these cases. The gains are small relative to the static grid and are not broad enough to justify deployment. The current conclusion is therefore negative for a zero-shot activation router: factors are useful for explaining conditional behavior and selecting follow-up ablations, but not for selecting a concrete activation and placement in a new cell.

\begin{table*}[t]
\centering
\caption{Top positive leave-one-seed-out factor-oracle diagnostics among array-ready Phase-4 cells. These are diagnostic upper bounds for regime sensitivity, not deployed activation-routing results.}
\label{tab:phase4_loso_oracle}
\resizebox{\textwidth}{!}{%
\begin{tabular}{lllccc}
\toprule
Dataset & H & factor & LOSO oracle MSE & best static MSE & gain \\
\midrule
Solar1 & 96 & max\_abs\_change & 0.174 & 0.178 & +2.0\% \\
Solar1 & 96 & volatility & 0.175 & 0.178 & +1.4\% \\
Solar1 & 96 & volatility\_shock & 0.175 & 0.178 & +1.3\% \\
Solar1 & 96 & trend\_consistency & 0.177 & 0.178 & +0.4\% \\
Solar1 & 96 & max\_drawdown & 0.177 & 0.178 & +0.2\% \\
Solar1 & 96 & max\_drawup & 0.177 & 0.178 & +0.2\% \\
Solar1 & 96 & range & 0.177 & 0.178 & +0.2\% \\
\bottomrule
\end{tabular}%
}
\end{table*}


\subsection{Architecture-transfer stress tests}

\Cref{tab:patchtst_transfer} and \Cref{tab:patchtst_output_tanh} test whether the Informer findings transfer to PatchTST. The answer is partial. In the FFN-linear stage, GELU remains best in 7 of 12 cells, while \texttt{softsign} and \texttt{tanh\_sin001} win selected ETTh2 and Solar cells. This shows that activation sensitivity does not disappear, but the Informer static ranking does not transfer directly. The output-placement result is stronger and negative: output \texttt{tanh} loses all 54 paired seed-run comparisons against the matching linear-head configuration.

\begin{table}[t]
\centering
\caption{PatchTST activation-transfer stress test. The Informer activation effect transfers only partially: non-GELU activations win selected ETTh2/Solar cells, while GELU remains best in 7 of 12 tested cells.}
\label{tab:patchtst_transfer}
\begin{tabular}{llrrlr}
\toprule
Model & Dataset & Horizon & GELU MSE & Best config & Gain \\
\midrule
PatchTST & ETTh2 & 168 & 0.325880 & \texttt{gelu\_ffn\_linear} & 0.00\% \\
PatchTST & ETTh2 & 336 & 0.328450 & \texttt{gelu\_ffn\_linear} & 0.00\% \\
PatchTST & ETTh2 & 720 & 0.380881 & \texttt{softsign\_ffn\_linear} & 0.67\% \\
PatchTST & ETTm2 & 96 & 0.165589 & \texttt{gelu\_ffn\_linear} & 0.00\% \\
PatchTST & ETTm2 & 288 & 0.262739 & \texttt{gelu\_ffn\_linear} & 0.00\% \\
PatchTST & ETTm2 & 672 & 0.355232 & \texttt{gelu\_ffn\_linear} & 0.00\% \\
PatchTST & Solar2 & 24 & 0.065450 & \texttt{tanh\_sin001\_ffn\_linear} & 0.42\% \\
PatchTST & Solar2 & 96 & 0.099666 & \texttt{softsign\_ffn\_linear} & 1.12\% \\
PatchTST & Solar3 & 24 & 0.279035 & \texttt{gelu\_ffn\_linear} & 0.00\% \\
PatchTST & Solar3 & 96 & 0.446409 & \texttt{gelu\_ffn\_linear} & 0.00\% \\
PatchTST & Solar5 & 24 & 0.114792 & \texttt{tanh\_sin001\_ffn\_linear} & 1.33\% \\
PatchTST & Solar5 & 96 & 0.176366 & \texttt{softsign\_ffn\_linear} & 1.01\% \\
\bottomrule
\end{tabular}
\end{table}

\begin{table}[t]
\centering
\caption{PatchTST output-tanh placement test. Output tanh loses every paired seed-run comparison against the matching linear-head configuration, indicating that the Informer output-bounding signal does not transfer to PatchTST.}
\label{tab:patchtst_output_tanh}
\begin{tabular}{lrrr}
\toprule
Base activation & Paired seed-runs & Win rate & Mean gain vs. linear \\
\midrule
\texttt{tanh} & 18 & 0.0\% & -5.28\% \\
\texttt{tanh\_sin001} & 18 & 0.0\% & -5.31\% \\
\texttt{gelu} & 18 & 0.0\% & -5.44\% \\
\bottomrule
\end{tabular}
\end{table}


\Cref{tab:itransformer_transfer} gives a second transfer stress test. In a one-seed ETTh2/ETTm2 FFN-only pilot, GELU remains best in five of six cells, and the only non-GELU best case is a narrow \texttt{tanh\_sin001} gain on ETTm2/96. A later three-cell placement probe gives a more conditional signal: output \texttt{tanh} helps on ETTh2/96 and ETTm2/720 but hurts on ETTm2/96. Because this placement probe is currently single-seed, we use it only as a direction for follow-up, not as a main statistical claim.

\begin{table}[t]
\centering
\caption{iTransformer activation-transfer stress test. In this one-seed pilot, GELU remains best in 5 of 6 ETTh2/ETTm2 cells; the only non-GELU win is a narrow \texttt{tanh\_sin001} gain on ETTm2/96.}
\label{tab:itransformer_transfer}
\begin{tabular}{llrrlr}
\toprule
Model & Dataset & Horizon & GELU MSE & Best config & Gain \\
\midrule
iTransformer & ETTh2 & 96 & 0.301598 & \texttt{gelu\_ffn\_linear} & 0.00\% \\
iTransformer & ETTh2 & 336 & 0.422894 & \texttt{gelu\_ffn\_linear} & 0.00\% \\
iTransformer & ETTh2 & 720 & 0.428986 & \texttt{gelu\_ffn\_linear} & 0.00\% \\
iTransformer & ETTm2 & 96 & 0.184226 & \texttt{tanh\_sin001\_ffn\_linear} & 0.19\% \\
iTransformer & ETTm2 & 336 & 0.313418 & \texttt{gelu\_ffn\_linear} & 0.00\% \\
iTransformer & ETTm2 & 720 & 0.414092 & \texttt{gelu\_ffn\_linear} & 0.00\% \\
\bottomrule
\end{tabular}
\end{table}


Together, the transfer experiments sharpen the story. The Informer full-grid claim is strong because it is complete and matched. The broader cross-architecture claim is weaker but still useful: activation behavior depends on architecture and placement, so an activation policy should not be promoted as a universal dataset-only rule.
